The quarterly estimates use Python \texttt{statsmodels.GLM} with the Poisson
family and iteratively reweighted least squares, followed by explicit
Poisson-score/Hessian cluster covariance calculations.  City all-zero groups
and singletons are removed before fitting.  A linear program checks whether a
design direction leaves positive-outcome predictors unchanged while weakly
reducing zero-outcome predictors, with strict reduction for at least one zero;
this is the separation criterion discussed by
\citet{correia2026separation}.  No strict separation direction is detected in
any of the five designs.  The complete covariance is transformed from the
fitted base-level/difference parameterization to the original-level
coefficients.

All specifications converge.  The common main support contains 1,301 positive
and 241 source-coded-zero outcomes: 1,542 rows in 111 cities over 28 calendar
quarters from 2017Q1 through 2023Q4, with 2--28 retained quarters per city and
a median of 11.  Every retained row has $W/W^{\mathrm{cal}}=1$; the minimum,
quartiles and maximum are all one.  Complete coverage was not imposed as an
additional sample threshold.  Main-model inference uses $t(110)$.

Before fixed-effect support exclusions, the four-quarter facility support has
2,229 rows in 273 cities, including 1,309 positive and 920 source-coded-zero
outcomes.  The main support removes 155 all-zero city groups (679 rows) and 7
singleton city groups (7 rows); one additional row fails the city/quarter
fixed-effect support condition.  The result is the 1,542 city-quarter support
in 111 cities.

Before fixed-effect support exclusions, the COVID support has 1,107 rows in 104
cities.  Fixed-effect support exclusions remove 17 rows: 13 rows meet the
city-all-zero condition and 5 meet the city-singleton condition, with one row
meeting both conditions.  Equivalently, the mutually exclusive count is 13
all-zero rows plus 4 additional singleton-only rows.  The resulting paired
support has 1,090 rows in 94 cities over 24
quarters and 25 province clusters; its inference uses $t(24)$.  No 2023 COVID
exposure is filled.  These exclusions define the reported estimands and are
separate from missing compatible weeks.
